\begin{abstract}
Attackers routinely evade Web Application Firewalls (WAFs) by mutating SQL injection (SQLi) payloads. Prior studies report high bypass rates, yet these aggregate figures often conflate heterogeneous operators and depend on search biases. Consequently, defenders need reproducible, family-resolved measurements tied to explicit labels.

This paper presents an offensive methodology that generates discrete SQLi variants through dialect seeds, deployment-specific augmentations, and learned transform rules. We issued black-box HTTP probes against a live AWS WAF-fronted application, collecting per-attempt JSONL logs. In our archived run of 5,000 variants, our operational classifier identified a 41.7\% bypass rate. However, family-level analysis reveals extreme heterogeneity: identity, case-mix, and null-byte cohorts achieved bypass rates near 99\%, while encoding and hex-style cohorts were entirely blocked. This indicates that policy stress concentrates in specific syntactic strategies.

As a secondary control, we applied repeated RFC percent-decoding and scored strings with a TF-IDF logistic-regression surrogate. On a 20-request pilot, accuracy remained at 42.9\% for both raw and canonicalized inputs, suggesting that naive decoding neither closes the live bypass gap nor rescues a simple offline model. All reported rates are conditional on the specific experimental configuration; the core contribution lies in the explicit methodology, decomposed statistics, and reproducible artifact.
\end{abstract}
